% ── WGS Tier 1 补充分析（插入 Section VI-C 末尾）────────────────
\paragraph{Tier-1 result under weighted flow time objective.}
To isolate the domain-specific effect of LPT near-optimality on the WGS
Tier-1 Beats-Best rate (0.6\% under makespan), we evaluate two additional
distributional scenarios (D11, D12) under a \emph{weighted flow time}
objective ($\sum_j w_j C_j$), where LPT is no longer a strong baseline.
D11 models a GPU-intensive task surge (20\%\,→\,65\% GPU-intensive tasks
with weight ratio 3$\times$), and D12 models an emergency batch influx
(5\%\,→\,40\% urgent tasks with $w_j = 5$).

Across 64 trials (8\,seeds $\times$ 2\,scenarios
$\times$ 4\,baselines), LAOSA achieves a
Beats-Best rate of
$\mathbf{73.4\%}$
and outperforms LPT specifically in
$\mathbf{100.0\%}$
of trials, with mean improvement
$8.4\% \pm 3.3\%$
over the degraded baseline.
This confirms that LAOSA's low Tier-1 Beats-Best under makespan
is domain-specific (LPT dominates the static baseline pool for
makespan minimisation), and not a limitation of the generative
adaptation mechanism itself.

% ── 补充小表（附在 Table IV 之后）────────────────────────────────
\begin{table}[h]
\centering
\caption{WGS Tier-1 under weighted flow time objective
(D11\,+\,D12, 64\,trials, 8\,seeds).
Beats Best: vs.\ best static; Beats LPT: vs.\ LPT specifically.
$\dagger$ LPT is near-optimal for makespan but not for WFT.}
\label{tab:tier1_wft}
\small
\begin{tabular}{@{}lccc@{}}
\toprule
Scenario & Beats Best & Beats LPT & Improvement (\%) \\
\midrule
  D11 (GPU surge, WFT)           & 62.5\% & 100.0\% & $8.2\pm2.4$ \\
  D12 (Emergency burst, WFT)     & 84.4\% & 100.0\% & $8.5\pm4.0$ \\
  LPT baseline (same scenarios)  & N/A & 0.0\%$^{\dagger}$ & $\approx 0$ \\
\midrule
\multicolumn{4}{l}{\footnotesize
  WGS Tier-1 (makespan, Table~\ref{tab:q1_bytype}): Beats Best = 0.6\% (2/320)}\\
\bottomrule
\end{tabular}
\end{table}